\section{Related work for the supporting studies}
\label{app:supporting-related}
\subsection{Earlier Koi-Koi agents}
Sato and Ikeda compared policy gradients with Neural Fitted Q Iteration using
268 features and simplified scoring \cite{sato2017koikoi}. Nagamine and Suzuki
used rollout rewards for fast-finishing and blocking strategies, with koi-koi
continuation disabled \cite{nagamine2021strategy}. Their proposed switching
between objectives motivates our comparison of learned selectors and fixed
card/koi roles, including continuation decisions. Takaoka et al. adapt UCT
choices to the score to keep games close and study enjoyment
\cite{takaoka2018fun}; their objective differs from match-win optimization.

\subsection{Adapting action scores and exploration}
Policy distillation trains a student on temperature-scaled teacher scores
\cite{rusu2016distillation}. We optimize the inherited actor directly with
match rewards. Discrete SAC derives categorical maximum-entropy learning
\cite{christodoulou2019discrete}, and target-entropy scheduling adapts exploration
\cite{xu2021entropy}. Our fixed score divisors test how published scores transfer
into match-trained policies. One-head and four-head calibration change both
sampling and policy-loss derivatives (Appendices~\ref{app:temperature}
and~\ref{sec:point-initialization}).

\subsection{Self-play and strategic interaction}
Neural Fictitious Self-Play learns best responses and average strategies
\cite{heinrich2016nfsp}; Policy-Space Response Oracles respond to population
mixtures selected through an empirical game \cite{lanctot2017psro}.
Counterfactual regret minimization connects local regret to approximate
equilibrium in two-player zero-sum games with perfect recall
\cite{zinkevich2007cfr}. Our current-policy self-play measures fixed-opponent
matchups; it supplies no exploitability analysis.

\subsection{Guided adaptation and policy selection}
Bayesian Policy Reuse selects policies from observed task signals
\cite{rosman2016reuse}; DRON learns opponent-dependent responses from opponent
observations \cite{he2016opponent}. Our public/masked comparison isolates the
recovered identity's contribution. Jump-Start Reinforcement Learning constructs
starting-state curricula with a guide \cite{uchendu2023jump}. Our fixed parent
prefixes similarly isolate adaptation after reaching the final round or
public identification.

Classification-based policy iteration \cite{lagoudakis2003classification},
adaptive rollout allocation \cite{dimitrakakis2008rollout} and reward-sensitive
classification \cite{langford2005classification} connect rollout outcomes to
policy learning. DAgger and AggreVaTe address the changing observation
distribution through interactive learning \cite{ross2011dagger,ross2014aggrevate}.
Our selector fits one fixed-source dataset before complete-match evaluation.
Proposals share hidden cards and future deals; this common-random-number
coupling can reduce comparison variance \cite{yadav2026common}.

Comanici and Precup learn switching through option termination
\cite{comanici2010switching}. Our decision-level selectors and fixed card/koi
hybrids test whether point-oriented and win-oriented experts provide useful
alternatives for match optimization.
